
\documentclass{article}
\usepackage{colortbl}
\usepackage{makecell}
\usepackage{multirow}
\usepackage{supertabular}

\begin{document}

\newcounter{utterance}

\centering \large Interaction Transcript for game `eqbench', experiment `eq\_bench', episode 61 with reflexion{-}llama{-}3.1{-}8b{-}instruct.
\vspace{24pt}

{ \footnotesize  \setcounter{utterance}{1}
\setlength{\tabcolsep}{0pt}
\begin{supertabular}{c@{$\;$}|p{.15\linewidth}@{}p{.15\linewidth}p{.15\linewidth}p{.15\linewidth}p{.15\linewidth}p{.15\linewidth}}
    \# & \multicolumn{2}{c}{Player} && \multicolumn{2}{c}{Game Master} \\
    \hline

    \theutterance \stepcounter{utterance}  
    & & & \multicolumn{4}{p{0.6\linewidth}}{
        \cellcolor[rgb]{0.9,0.9,0.9}{
            \makecell[{{p{\linewidth}}}]{
                \texttt{\tiny{[P1$\langle$GM]}}
                \texttt{You have attempted this same task before and did not succeed. These are the plans you wrote after those attempts.} \\
\texttt{Trial \#1: New plan: Upon re{-}reading the dialogue, I realized that I focused too much on the surface{-}level conversation and didn't consider the underlying dynamics between Marcus and Julia. I should have paid more attention to the power struggle and the themes of truth, danger, and fear. Next time, I will analyze the dialogue more deeply, looking for subtle cues and underlying emotions. I will also consider the context of the situation (being stuck in an elevator) and how it affects the characters' emotional states. Specifically, I will examine Julia's responses more closely, as they seem to be the key to understanding her emotional state. I will also consider the tone and language used by Marcus, as his warnings and threats may have a significant impact on Julia's emotions.} \\
\\ 
\texttt{Trial \#2: New plan: Upon re{-}reading my previous attempt, I realized that I scored the emotions based on my initial impression of the dialogue, but I didn't consider the progression of the conversation and the character development. I should have paid more attention to the way Julia responds to Marcus's warnings and threats, and how her tone and language change throughout the dialogue. Specifically, I should have noticed that Julia's responses become more confident and assertive as the conversation progresses, which suggests that she is not feeling fear, but rather defiance. I will also consider the fact that Julia's emotional state is not static, but rather dynamic, and that her emotions may shift in response to the situation and the other character's actions. Next time, I will analyze the dialogue more closely, looking for subtle cues and underlying emotions, and I will consider the context of the situation and the character development to make a more accurate prediction of Julia's emotional state.} \\
\\ 
\\ 
\\ 
\texttt{Your task is to predict the likely emotional responses of a character in this dialogue:} \\
\\ 
\texttt{Marcus: Stuck in an elevator. Just my luck.} \\
\texttt{Julia: Could be worse. Could be stuck in a room with no windows.} \\
\texttt{Marcus: You're an optimist.} \\
\texttt{Julia: Realist. I just choose to see the glass half full.} \\
\texttt{Marcus: And what if the glass is poisoned?} \\
\texttt{Julia: Then it's a good thing I'm thirsty.} \\
\texttt{Marcus: You always this reckless?} \\
\texttt{Julia: Only when it matters.} \\
\texttt{Marcus: And what matters to you?} \\
\texttt{Julia: The truth.} \\
\texttt{Marcus: Truth is a dangerous game.} \\
\texttt{Julia: Only for those with something to hide.} \\
\texttt{Marcus: Or those who seek it.} \\
\texttt{Julia: I'm not afraid.} \\
\texttt{Marcus: You should be.} \\
\texttt{Julia: Is that a threat?} \\
\texttt{Marcus: A warning.} \\
\texttt{{[}End dialogue{]}} \\
\\ 
\texttt{At the end of this dialogue, Julia would feel...} \\
\texttt{Fear} \\
\texttt{Defiance} \\
\texttt{Indifference} \\
\texttt{Intrigue} \\
\\ 
\texttt{Give each of these possible emotions a score from 0{-}10 for the relative intensity that they are likely to be feeling each.} \\
\\ 
\texttt{You must output in the following format, including headings (of course, you should give your own scores), with no additional commentary:} \\
\\ 
\texttt{Fear: <score>} \\
\texttt{Defiance: <score>} \\
\texttt{Indifference: <score>} \\
\texttt{Intrigue: <score>} \\
\\ 
\\ 
\texttt{{[}End of answer{]}} \\
\\ 
\texttt{Remember: zero is a valid score, meaning they are likely not feeling that emotion. You must score at least one emotion > 0.} \\
\\ 
\texttt{Your answer:} \\
            }
        }
    }
     \\ \\

    \theutterance \stepcounter{utterance}  
    & \multicolumn{4}{p{0.6\linewidth}}{
        \cellcolor[rgb]{0.561,0.737,0.561}{
            \makecell[{{p{\linewidth}}}]{
                \texttt{\tiny{[P1$\rangle$GM]}}
                \texttt{Fear: 2} \\
\texttt{Defiance: 8} \\
\texttt{Indifference: 0} \\
\texttt{Intrigue: 0} \\
            }
        }
    }
    & & \\ \\

    \theutterance \stepcounter{utterance}  
    & & & \multicolumn{2}{p{0.3\linewidth}}{
        \cellcolor[rgb]{0.9,0.9,0.9}{
            \makecell[{{p{\linewidth}}}]{
                \texttt{\tiny{[GM$|$GM]}}
                \texttt{Fear: 2} \\
\texttt{Defiance: 8} \\
\texttt{Indifference: 0} \\
\texttt{Intrigue: 0} \\
            }
        }
    }
    & & \\ \\

    \theutterance \stepcounter{utterance}  
    & & & \multicolumn{2}{p{0.3\linewidth}}{
        \cellcolor[rgb]{0.9,0.9,0.9}{
            \makecell[{{p{\linewidth}}}]{
                \texttt{\tiny{[GM$|$GM]}}
                \texttt{Fear: 3} \\
\texttt{Defiance: 7} \\
\texttt{Indifference: 0} \\
\texttt{Intrigue: 6} \\
            }
        }
    }
    & & \\ \\

    \theutterance \stepcounter{utterance}  
    & & & \multicolumn{2}{p{0.3\linewidth}}{
        \cellcolor[rgb]{0.9,0.9,0.9}{
            \makecell[{{p{\linewidth}}}]{
                \texttt{\tiny{[GM$|$GM]}}
                \texttt{game\_result = LOSE} \\
            }
        }
    }
    & & \\ \\

    \theutterance \stepcounter{utterance}  
    & & \multicolumn{4}{p{0.6\linewidth}}{
        \cellcolor[rgb]{0.95,0.95,0.95}{
            \makecell[{{p{\linewidth}}}]{
                \texttt{\tiny{[P1$\leftrightarrow$P]}}
                \texttt{Post{-}episode Reflexion (not sent to GM):} \\
\texttt{New plan: Upon re{-}reading my previous attempts, I realized that I focused too much on the surface{-}level conversation and didn't consider the underlying dynamics between Marcus and Julia. I should have paid more attention to the power struggle and the themes of truth, danger, and fear. However, I also made the mistake of not considering the context of the situation (being stuck in an elevator) and how it affects the characters' emotional states. } \\
\\ 
\texttt{In my previous attempts, I scored the emotions based on my initial impression of the dialogue, but I didn't consider the progression of the conversation and the character development. I should have paid more attention to the way Julia responds to Marcus's warnings and threats, and how her tone and language change throughout the dialogue.} \\
\\ 
\texttt{To improve my prediction, I will analyze the dialogue more closely, looking for subtle cues and underlying emotions. I will consider the context of the situation and the character development to make a more accurate prediction of Julia's emotional state. Specifically, I will:} \\
\\ 
\texttt{1. Examine Julia's responses more closely, as they seem to be the key to understanding her emotional state.} \\
\texttt{2. Consider the tone and language used by Marcus, as his warnings and threats may have a significant impact on Julia's emotions.} \\
\texttt{3. Analyze the progression of the conversation and the character development to understand how Julia's emotional state changes throughout the dialogue.} \\
\texttt{4. Consider the fact that Julia's emotional state is not static, but rather dynamic, and that her emotions may shift in response to the situation and the other character's actions.} \\
\\ 
\texttt{By following this plan, I hope to make a more accurate prediction of Julia's emotional state and improve my overall performance in this task.} \\
            }
        }
    }
    & \\ \\

\end{supertabular}
}

\end{document}
